% Part: sets-functions-relations
% Chapter: sets
% Section: subsets

\documentclass[../../../include/open-logic-section]{subfiles}

\begin{document}

\olfileid{sfr}{set}{sub}
\olsection{ذیلی سیٹ تے ذیلی سیٹاں دے سیٹ}

\begin{explain}
سانوں اکثر سیٹاں دا آپس وچ ٹاکرا کرن دی لوڑ پئے گی۔ ٹاکرے دی اک صاف صورت
ایہ اے: \emph{اک سیٹ وچ جیہڑی وی شے اے، اوہ دوجے وچ وی اے}۔ ایہ صورت اینی
اہم اے کہ ایس لئی کجھ نویاں علامتاں متعارف کرائیاں جان۔
\end{explain}

\begin{defn}[ذیلی سیٹ]
جے سیٹ $A$ دا ہر !!{element}، $B$ دا وی !!a{element} ہووے، تاں اسی آکھدے آں
کہ $A$، $B$ دا \emph{ذیلی سیٹ} اے، تے $A \subseteq B$ لکھدے آں۔ جے
$A$، $B$ دا ذیلی سیٹ نہ ہووے، تاں اسی $A \not\subseteq B$ لکھدے آں۔
جے $A \subseteq B$ ہووے پر $A \neq B$ ہووے، تاں اسی $A \subsetneq B$
لکھدے آں تے آکھدے آں کہ $A$، $B$ دا \emph{حقیقی ذیلی سیٹ} اے۔
\end{defn}

\begin{ex}
ہر سیٹ اپنے آپ دا ذیلی سیٹ اے، تے $\emptyset$ ہر سیٹ دا ذیلی سیٹ اے۔ قدرتی
جفت اعداد دا سیٹ قدرتی اعداد دے سیٹ دا ذیلی سیٹ اے۔ ایس طرح
$\{ a, b \} \subseteq \{ a, b, c \}$۔ پر $\{ a, b, e \}$،
$\{ a, b, c \}$ دا ذیلی سیٹ نہیں اے۔
\end{ex}

\begin{ex}
عدد $2$ صحیح اعداد دے سیٹ دا !!{element} اے، جد کہ جفت اعداد دا سیٹ صحیح
اعداد دے سیٹ دا ذیلی سیٹ اے۔ پر ایہ وی ہو سکدا اے کہ اک سیٹ کِسے ہور سیٹ
دا !!a{element} وی ہووے تے اوہدا ذیلی سیٹ وی، یعنی \emph{دونوں} گلاں ہون؛
مثال لئی، $\{0\} \in \{0, \{0\}\}$ تے نال ہی
$\{0\} \subseteq \{0, \{0\}\}$۔
\end{ex}

عنصراں نال تعیّن دا اصول سیٹاں دے اکو ہون دی کسوٹی دیندا اے: $A = B$ اُدوں
تے صرف اُدوں ہوندا اے جدوں $A$ دا ہر !!{element}، $B$ دا وی !!a{element}
ہووے تے ایس توں اُلٹ وی۔ ``ذیلی سیٹ'' دی تعریف $A \subseteq B$ نوں عین ایس
کسوٹی دے پہلے ادھ دے طور تے دسدی اے: $A$ دا ہر !!{element}، $B$ دا وی
!!a{element} اے۔ بے شک، جے اسی $A$ تے $B$ نوں آپس وچ وٹا دیئیے، تاں وی
تعریف لاگو ہوندی اے: یعنی $B \subseteq A$ اُدوں تے صرف اُدوں جدوں $B$ دا
ہر !!{element}، $A$ دا وی !!a{element} ہووے۔ تے ایہ عین عنصراں نال تعیّن
دے اصول دا ``ایس توں اُلٹ وی'' والا حصہ اے۔ ہور لفظاں وچ، ایس اصول توں
نکلدا اے کہ سیٹ اُدوں تے صرف اُدوں برابر نیں جدوں اوہ اک دوجے دے ذیلی سیٹ ہون۔

\begin{prop}
$A = B$ اُدوں تے صرف اُدوں ہوندا اے جدوں $A \subseteq B$ تے
$B \subseteq A$، دونوں ہون۔
\end{prop}

ہُن کجھ ہور مددگار علامتاں متعارف کراؤن دا وی چنگا موقع اے۔ ایہ تعریف
کردیاں کہ $A$، $B$ دا ذیلی سیٹ کدوں اے، اسی آکھیا سی کہ ``$A$ دا ہر
!!{element} \dots اے''، تے ``$\dots$'' دی تھاں ``$B$ دا !!a{element}''
رکھیا سی۔ پر عبارت دی ایہ \emph{بنت} اینی عام اے کہ ایس لئی کجھ رسمی
علامتاں متعارف کرانا مددگار ہووے گا۔

\begin{defn}\ollabel{forallxina}
$(\forall x \in A)\phi$، $\forall x(x \in A \lif
\phi)$ دی مختصر لکھت اے۔ ایس طرح، $(\exists x \in A)\phi$،
$\exists x(x \in A \land \phi)$ دی مختصر لکھت اے۔
\end{defn}

ایہ علامتاں ورت کے، اسی آکھ سکدے آں کہ $A \subseteq B$ اُدوں تے صرف اُدوں
ہوندا اے جدوں $(\forall x \in A)x \in B$ ہووے۔

ہُن اسی سیٹ دی اک خاص قسم ول آندے آں: کِسے دتے ہوئے سیٹ دے سارے ذیلی
سیٹاں دا سیٹ۔

\begin{defn}[ذیلی سیٹاں دا سیٹ]
سیٹ~$A$ دے سارے ذیلی سیٹاں توں بنے سیٹ نوں $A$ دا \emph{ذیلی سیٹاں دا سیٹ}
آکھدے نیں تے $\Pow{A}$ لکھدے نیں۔
  \[
    \Pow{A} = \Setabs{B}{B \subseteq A}
  \]
\end{defn}

\begin{ex}
$\{ a, b, c \}$ دے سارے ممکن ذیلی سیٹ کِہڑے نیں؟ اوہ ایہ نیں:
$\emptyset$، $\{a \}$، $\{b\}$، $\{c\}$، $\{a, b\}$، $\{a, c\}$،
$\{b, c\}$، $\{a, b, c\}$۔ ایہناں سارے ذیلی سیٹاں دا سیٹ
$\Pow{\{a,b,c\}}$ اے:
\[
\Pow{\{ a, b, c \}} = \{\emptyset, \{a \}, \{b\}, \{c\}, \{a, b\},
\{b, c\}, \{a, c\}, \{a, b, c\}\}
\]
\end{ex}

\begin{prob}
$\{a, b, c, d\}$ دے سارے ذیلی سیٹ لکھو۔
\end{prob}

\begin{prob}
وکھاؤ کہ جے $A$ دے $n$ !!{element}s ہون، تاں $\Pow{A}$ دے $2^n$
!!{element}s ہون گے۔
\end{prob}

\end{document}
